\section{Dataset: Multi-Reference 2D Sketch Constraint Annotations}
\label{sec:dataset}

We construct a 2D CAD sketch dataset to evaluate constraint quality with respect to design intent.
Unlike existing CAD datasets, it provides multiple independent expert annotations for the same sketch geometry, letting us study the ambiguity of design-intent-aware constraints.
Each of 292 sketches is annotated by three of six CAD experts into a fully constrained sketch, yielding 876 human annotations.

\subsection{Collection and Annotation}

\textbf{Source.}
We draw sketch geometries from Fusion 360 Gallery~\cite{willis2020fusion} rather than an Onshape-derived corpus (SketchGraphs~\cite{seff2020sketchgraphs}, ABC~\cite{Koch_2019_CVPR}, DeepCAD~\cite{Wu_2021_ICCV}) to limit source overlap with the training data of existing sketch-constraint models (App.~\ref{sec:appendix_dataset_source}).
We extract individual sketches from the original 3D CAD models and convert them into the FreeCAD format, since the same FreeCAD constraint solver is later used for solver-based validation and reinforcement-learning feedback (\S~\ref{sec:experiments},\ref{sec:rl_reward}).
Following Vitruvion~\cite{seff2021vitruvion}, we keep four primitive types (\texttt{Point}, \texttt{Line}, \texttt{Circle}, \texttt{Arc}); after discarding unsupported or degenerate sketches, $1{,}448$ valid sketches remain, most under-constrained (mean DoF $\approx 14.3$).
We sample $300$ while preserving the DoF distribution, avoiding a bias toward trivially or excessively constrained sketches; the $292$ sketches whose three annotations were completed and passed quality screening form the benchmark.

\textbf{Annotation.}
Each sketch is assigned to a varying trio of the six experts, yielding three independent annotations for the same input geometry.
The annotation protocol asks annotators to preserve the input geometry while adding design-intent-aware constraints, following the detailed guidelines in App.~\ref{sec:appendix_dataset_guidelines}.
In particular, annotators are prohibited from referencing the global axes or origin, which would obscure sketch-intrinsic design intent and reduce reusability.
This restriction intentionally leaves the two global translation degrees of freedom, so annotators add a single terminal \texttt{Block} constraint only to satisfy FreeCAD's fully constrained solver state (DoF $=0$); we treat this anchor as a technical artifact and exclude it from design-intent analysis.
For quality control, annotators may delete and replace existing constraints when doing so better captures design intent, but may not delete, redraw, or reorder the original primitives.
Adding up to two construction primitives is allowed only when necessary to express structural relations such as mirror axes, symmetry centers, or circular arrangements.

\textbf{Analysis schema.}
Although annotation is authored in FreeCAD, we report statistics in a SketchGraphs-style schema, retaining FreeCAD-derived labels such as \texttt{DistanceX}/\texttt{DistanceY} when constraint attributes distinguish them. The terminal \texttt{Block} is excluded as an artifact (App.~\ref{sec:appendix_dataset_schema}).

\subsection{Statistics and Ambiguity}

\begin{figure}[t]
    \centering
    \includegraphics[width=0.92\linewidth]{figs/ch3_constraint_distribution.pdf}
    \caption{Constraint-type distribution of the annotations (SketchGraphs schema), as a share of all constraints, grouped into relational (top) and dimensional (bottom) families.}
    \label{fig:constraint_dist}
\end{figure}

\textbf{Composition.}
Every annotation is fully constrained (DoF $=0$), with on average $24.1$ constraints over $8.6$ primitives.
As Fig.~\ref{fig:constraint_dist} shows, the annotations are relation-rich: relational constraints account for $75.2\%$ of all constraints.
Extended per-type, primitive, and per-annotator statistics are reported in App.~\ref{sec:appendix_dataset_statistics}.

\textbf{Ambiguity.}
The multi-reference design exposes constraint-selection ambiguity directly: for fixed geometry, experts choose different but valid strategies (e.g., equality vs.\ explicit dimensions).
As a result, disagreement is noticeable at the constraint-token level: the mean pairwise constraint-token F1 across the three expert annotations is $0.61$.
This shows why single-reference evaluation is problematic: a prediction that differs from one human annotation may still express another plausible design intent.
Multiple references therefore help distinguish invalid or poorly constrained predictions from valid alternatives chosen by different experts.
At the same time, token-level disagreement does not necessarily imply different redesign behavior; many token-divergent annotations induce similar local edit motions, motivating the behavior-based metrics introduced next.
